The use of the DecoupleR model (\cite{badia-i-mompel_decoupler_2022, badia-i-mompel_transcription_nodate}) raises some suspicion.
The model itself assumes that TF act individually to regulate a set of target genes.
A prior knowledge matrix, in our case based on CollecTRI (\cite{muller-dott_expanding_2023}),
describes which TF are known to up- or downregulate which genes.
The gene expressions of all the TF's target genes are then use to ``reverse engineer'' (estimate) the TF activity.
DecoupleR's output is taken as an undisputed truth onwards.
It is unclear whether combinatoric effects can be observed from output of a model that does not assume combinatoric behavior.

Furthermore, the rules we have found may fail to capture real-life phenomena
and instead reflect information from the prior knowledge matrix.
Due to the construction of the DecoupleR model, TFs that target similar genes will have similar activity estimates.
This may partly explain the high cross-correlations observed in \autoref{fig:corr_activity}.
As a result, even if TF activity differs in reality,
the DecoupleR model can make TF activities hard to distinguish.
This, in turn, can make finding meaningful rules difficult.
Instead, the rules may represent the prior knowledge matrix, meaning nothing of interest is learned.

As a curiosity, we use the threshold model (\autoref{ch:threshold-model}) on the raw gene expression data instead.
We do not use raw read counts but normalized TPM counts, meant to be comparable across samples.
While this is not expected to yield a result with a realistic interpretation,
it can possibly give insight in how the model behaves.

\autoref{fig:tpm_accuracy_mw_0.svg} shows the results, in the same way as \autoref{fig:threshold_accuracy_mw_0.svg} does.
The interpretable results for $C=1$ and $D=1,2,3$ are given in \autoref{tab:tpm_interpretable_results_C_1}.
The interpretable results for $C=2$ are given in \autoref{tab:tpm_interpretable_results_C_2}.

It follows that the accuracy achieved is higher compared to the experiments with TF activity estimates.
Furthermore, qualitatively different results are observed.
This indicates that gene expression is indeed not equivalent to TF activity.

\begin{figure}[h]
    \centering
    \includesvg[width=0.7\textwidth]{gen/PDL1ex3/PDL1ex3_fig_accuracy_mw_0.svg}
    \caption{Accuracy of results for the full Broad dataset, using gene expression as input, for $C=1,2,3,4,5$, $D=1,2,3,4$ and $Q_y=0.5,0.6,0.7,0.8,0.9.$}
    \label{fig:tpm_accuracy_mw_0.svg}
\end{figure}

\begin{table}[h]
    \centering
    \bwtable{gen/PDL1ex3/PDL1ex3_tab_interpretable_mw_0_C_1}
    \caption{Results for the full Broad dataset, using gene expression as input, for $C=1$ and small rules ($D=1,2,3$).}
    \label{tab:tpm_interpretable_results_C_1}
\end{table}

\begin{table}[h]
    \bwtable{gen/PDL1ex3/PDL1ex3_tab_interpretable_mw_0_C_2}
    \caption{Results for the full Broad dataset, using gene expression input, for $C=2$ and small rules ($D=1,2,3$).}
    \label{tab:tpm_interpretable_results_C_2}
\end{table}

Train accuracy is not the full story.
It remains unclear whether the results are biologically meaningful.
To study whether these results are meaningful in a mathematical sense,
we perform the exclusion experiment on this dataset too.
If there is a substantial drop in performance upon exclusion of a TF or a combination of TFs,
then this suggests that there do exist rules for the gene expression dataset that are more important than others.

The TF exclusion tables for $Q_y=0.5,0.6$ and $D=2$, the same parameters used earlier,
can be found in \autoref{tab:exclusion_tpm}.
Accuracy decreases remain small (between $0$ and $2$ percentage points), with two notable exceptions:
FOSL1 and IRF1.
In particular, for $Q_y=0.6$, excluding FOSL1 and IRF1 results in a substantial performance drop.
In the case $D=1$, exclusion of FOSL1 causes maximum attainable train accuracy to drop $9$ percentage points.
Then removing IRF1 has little effect.
More notably, for $D=2$, excluding FOSL1 and IRF1 causes an $8$ percentage point drop in the left branch and a
$9$ percentage point drop in the right branch, which we consider a substantial decrease.

This suggests that FOSL1 is an important TF with respect to this dataset, and that FOSL1 and IRF1 are an important
combination.
Furthermore, it shows that an exclusion experiment can yield new insights into TF combinations.
By extension, it strengthens the idea that TF activity estimates, as used in previous chapters,
are not suitable for inferring rulesets.

\clearpage
\begin{landscape}
    \vspace*{\fill}
    \begin{center}
        \captionof{table}{TF exclusion table as described in \autoref{sec:exclusion} for raw gene expression data,
            as described in \autoref{sec:gene_expression}.
        \\ Case $Q_y=0.5$ (top) and $Q_y=0.6$ (bottom), $D=1$ (left) and $D=2$ (right).}
        \vspace{1.5em}

        \label{tab:exclusion_tpm}
        \begin{minipage}{0.1\linewidth}
            \centering
            \begin{tabular}{||c|}\hline
\multicolumn{1}{|>{\columncolor[HTML]{B8E6B8}}c|}{Baseline: } \\
\multicolumn{1}{|>{\columncolor[HTML]{B8E6B8}}c|}{0.71} \\ \hline
\multicolumn{1}{|>{\columncolor[HTML]{F0BBB8}}c|}{-FOSL1:} \\
\multicolumn{1}{|>{\columncolor[HTML]{F0BBB8}}c|}{0.65} \\ \hline
\multicolumn{1}{|>{\columncolor[HTML]{F1BAB8}}c|}{-IRF1:} \\
\multicolumn{1}{|>{\columncolor[HTML]{F1BAB8}}c|}{0.65} \\ \hline
\multicolumn{1}{|>{\columncolor[HTML]{F5B8B8}}c|}{-STAT1:} \\
\multicolumn{1}{|>{\columncolor[HTML]{F5B8B8}}c|}{0.64} \\ \hline
\end{tabular}
        \end{minipage}
        \hfill
        \begin{minipage}{0.85\linewidth}
            \centering
            \begin{tabular}{||c|c|c|c|c|c|c|c|}\hline
\multicolumn{8}{|>{\columncolor[HTML]{B8E6B8}}c|}{Baseline: } \\
\multicolumn{8}{|>{\columncolor[HTML]{B8E6B8}}c|}{0.73} \\ \hline
\multicolumn{4}{|>{\columncolor[HTML]{DACCB8}}c|}{-FOSL1:} & \multicolumn{4}{|>{\columncolor[HTML]{B9E5B8}}c|}{-STAT4:} \\
\multicolumn{4}{|>{\columncolor[HTML]{DACCB8}}c|}{0.70} & \multicolumn{4}{|>{\columncolor[HTML]{B9E5B8}}c|}{0.73} \\ \hline
\multicolumn{2}{|>{\columncolor[HTML]{E7C2B8}}c|}{-IRF1:} & \multicolumn{2}{|>{\columncolor[HTML]{DEC8B8}}c|}{-NFATC1:} & \multicolumn{2}{|>{\columncolor[HTML]{DACCB8}}c|}{-FOSL1:} & \multicolumn{2}{|>{\columncolor[HTML]{BCE2B8}}c|}{-NFATC1:} \\
\multicolumn{2}{|>{\columncolor[HTML]{E7C2B8}}c|}{0.69} & \multicolumn{2}{|>{\columncolor[HTML]{DEC8B8}}c|}{0.70} & \multicolumn{2}{|>{\columncolor[HTML]{DACCB8}}c|}{0.70} & \multicolumn{2}{|>{\columncolor[HTML]{BCE2B8}}c|}{0.73} \\ \hline
\multicolumn{1}{|>{\columncolor[HTML]{F5B8B8}}c|}{-STAT1:} & \multicolumn{1}{|>{\columncolor[HTML]{F0BBB8}}c|}{-NFATC1:} & \multicolumn{1}{|>{\columncolor[HTML]{F0BBB8}}c|}{-IRF1:} & \multicolumn{1}{|>{\columncolor[HTML]{E0C7B8}}c|}{-SPI1:} & \multicolumn{1}{|>{\columncolor[HTML]{E7C2B8}}c|}{-IRF1:} & \multicolumn{1}{|>{\columncolor[HTML]{DEC8B8}}c|}{-NFATC1:} & \multicolumn{1}{|>{\columncolor[HTML]{DEC8B8}}c|}{-FOSL1:} & \multicolumn{1}{|>{\columncolor[HTML]{C3DDB8}}c|}{-IRF1:} \\
\multicolumn{1}{|>{\columncolor[HTML]{F5B8B8}}c|}{0.68} & \multicolumn{1}{|>{\columncolor[HTML]{F0BBB8}}c|}{0.68} & \multicolumn{1}{|>{\columncolor[HTML]{F0BBB8}}c|}{0.68} & \multicolumn{1}{|>{\columncolor[HTML]{E0C7B8}}c|}{0.70} & \multicolumn{1}{|>{\columncolor[HTML]{E7C2B8}}c|}{0.69} & \multicolumn{1}{|>{\columncolor[HTML]{DEC8B8}}c|}{0.70} & \multicolumn{1}{|>{\columncolor[HTML]{DEC8B8}}c|}{0.70} & \multicolumn{1}{|>{\columncolor[HTML]{C3DDB8}}c|}{0.72} \\ \hline
\end{tabular}
        \end{minipage}

        \vspace{1.5em}

        \begin{minipage}{0.1\linewidth}
            \centering
            \begin{tabular}{||c|}\hline
\multicolumn{1}{|>{\columncolor[HTML]{B8E6B8}}c|}{Baseline: } \\
\multicolumn{1}{|>{\columncolor[HTML]{B8E6B8}}c|}{0.74} \\ \hline
\multicolumn{1}{|>{\columncolor[HTML]{E9C0B8}}c|}{-FOSL1:} \\
\multicolumn{1}{|>{\columncolor[HTML]{E9C0B8}}c|}{0.65} \\ \hline
\multicolumn{1}{|>{\columncolor[HTML]{F0BBB8}}c|}{-IRF1:} \\
\multicolumn{1}{|>{\columncolor[HTML]{F0BBB8}}c|}{0.64} \\ \hline
\multicolumn{1}{|>{\columncolor[HTML]{F5B8B8}}c|}{-STAT4:} \\
\multicolumn{1}{|>{\columncolor[HTML]{F5B8B8}}c|}{0.63} \\ \hline
\end{tabular}
        \end{minipage}
        \hfill
        \begin{minipage}{0.85\linewidth}
            \centering
            \begin{tabular}{||c|c|c|c|c|c|c|c|}\hline
\multicolumn{8}{|>{\columncolor[HTML]{B8E6B8}}c|}{Baseline: } \\
\multicolumn{8}{|>{\columncolor[HTML]{B8E6B8}}c|}{0.76} \\ \hline
\multicolumn{4}{|>{\columncolor[HTML]{D7CEB8}}c|}{-FOSL1:} & \multicolumn{4}{|>{\columncolor[HTML]{BEE1B8}}c|}{-NFATC1:} \\
\multicolumn{4}{|>{\columncolor[HTML]{D7CEB8}}c|}{0.71} & \multicolumn{4}{|>{\columncolor[HTML]{BEE1B8}}c|}{0.75} \\ \hline
\multicolumn{2}{|>{\columncolor[HTML]{E7C1B8}}c|}{-IRF1:} & \multicolumn{2}{|>{\columncolor[HTML]{D8CDB8}}c|}{-NFATC1:} & \multicolumn{2}{|>{\columncolor[HTML]{D8CDB8}}c|}{-FOSL1:} & \multicolumn{2}{|>{\columncolor[HTML]{BFE0B8}}c|}{-IRF1:} \\
\multicolumn{2}{|>{\columncolor[HTML]{E7C1B8}}c|}{0.68} & \multicolumn{2}{|>{\columncolor[HTML]{D8CDB8}}c|}{0.71} & \multicolumn{2}{|>{\columncolor[HTML]{D8CDB8}}c|}{0.71} & \multicolumn{2}{|>{\columncolor[HTML]{BFE0B8}}c|}{0.75} \\ \hline
\multicolumn{1}{|>{\columncolor[HTML]{E9C0B8}}c|}{-STAT1:} & \multicolumn{1}{|>{\columncolor[HTML]{F5B8B8}}c|}{-NFATC1:} & \multicolumn{1}{|>{\columncolor[HTML]{F5B8B8}}c|}{-IRF1:} & \multicolumn{1}{|>{\columncolor[HTML]{DCCAB8}}c|}{-MYC:} & \multicolumn{1}{|>{\columncolor[HTML]{F5B8B8}}c|}{-IRF1:} & \multicolumn{1}{|>{\columncolor[HTML]{DCCAB8}}c|}{-MYC:} & \multicolumn{1}{|>{\columncolor[HTML]{F5B8B8}}c|}{-FOSL1:} & \multicolumn{1}{|>{\columncolor[HTML]{C3DDB8}}c|}{-STAT4:} \\
\multicolumn{1}{|>{\columncolor[HTML]{E9C0B8}}c|}{0.68} & \multicolumn{1}{|>{\columncolor[HTML]{F5B8B8}}c|}{0.66} & \multicolumn{1}{|>{\columncolor[HTML]{F5B8B8}}c|}{0.66} & \multicolumn{1}{|>{\columncolor[HTML]{DCCAB8}}c|}{0.70} & \multicolumn{1}{|>{\columncolor[HTML]{F5B8B8}}c|}{0.66} & \multicolumn{1}{|>{\columncolor[HTML]{DCCAB8}}c|}{0.70} & \multicolumn{1}{|>{\columncolor[HTML]{F5B8B8}}c|}{0.66} & \multicolumn{1}{|>{\columncolor[HTML]{C3DDB8}}c|}{0.74} \\ \hline
\end{tabular}
        \end{minipage}
    \end{center}
    \vspace*{\fill}
\end{landscape}